% GENERATED FILE. Edit canonical lessons, then run npm run generate:books.
% canonical-source-hash: fnv1a64:cdc7d98f8631cf84
% canonical-lessons: FR-C357-stage, FR-C357-stagiaire, FR-C357-apprenti, FR-C357-apprentissage, FR-C357-employeur, FR-R357-revision-350-357

\chapter{Le premier poste}
\label{ch:fr-le-premier-poste}
\hlchaptermodality{357}

\begin{chapteropening}
\textbf{By the end of this chapter:} \emph{I can talk about a first job: an internship, an intern, an apprentice, an apprenticeship and the employer.}

Complete the last lesson of chapter 357: I can talk about a first job: an internship, an intern, an apprentice, an apprenticeship and the employer.
\end{chapteropening}

\section[le stage]{le stage — the internship, the work placement; a training course}

\label{lesson:FR-C357-stage}



\begin{quote}
Before the new one: say the French for to hire, then the French for to recruit.
\end{quote}

\subsection*{You'll want to know: le stage}
\textbf{le stage} — \textquotedblleft{}the internship, the work placement; a training course\textquotedblright{}.

\begin{quote}
Cet été, je fais un stage dans un hôtel.
\end{quote}

\begin{quote}
\emph{This summer, I am doing an internship in a hotel.}
\end{quote}

A false friend: not a theatre stage. \emph{Un stage de cuisine} — a cooking course.

\begin{cousinweb}
\emph{stage} comes from Medieval Latin \emph{stagium} \textquotedblleft{}a period of residence\textquotedblright{} (EtymDB), usually linked to Latin \emph{stāre} \textquotedblleft{}to stand\textquotedblright{} (Lewis \& Short).
\end{cousinweb}

\begin{grammarlens}[title={What you've built}]
One more word for telling, explaining and describing at B1.
\end{grammarlens}

\subsection*{Your turn}
\begin{itemize}
\raggedright
  \item \emph{Say it:} \emph{le stage}
  \item \emph{Say it:} \emph{le stage}, once more
  \item \emph{Say it:} a sentence of your own with \emph{le stage}
  \item \emph{Recall:} say the French for to hire, then the French for to recruit, then say \emph{le stage} again
\end{itemize}

\begin{tcolorbox}[breakable,colback=teal!4,colframe=teal!35!black,title={Before you move on}]
What does le stage mean? (\textquotedblleft{}The internship, the work placement; a training course\textquotedblright{}.) Say it once more.
\end{tcolorbox}
\section[le stagiaire]{le stagiaire — the intern, the trainee}

\label{lesson:FR-C357-stagiaire}



\begin{quote}
Before the new one: say the French for to recruit, then the French for the internship.
\end{quote}

\subsection*{You'll want to know: le stagiaire}
\textbf{le stagiaire} — \textquotedblleft{}the intern, the trainee\textquotedblright{}.

\begin{quote}
Le stagiaire travaille avec moi depuis lundi.
\end{quote}

\begin{quote}
\emph{The intern has been working with me since Monday.}
\end{quote}

The same form for a man or a woman: \emph{une stagiaire}.

\begin{cousinweb}
\emph{stagiaire} is built from \emph{stage} and \emph{-aire}. \emph{stage} comes from Medieval Latin \emph{stagium} \textquotedblleft{}a period of residence\textquotedblright{} (EtymDB), usually linked to Latin \emph{stāre} \textquotedblleft{}to stand\textquotedblright{}.
\end{cousinweb}

\begin{grammarlens}[title={What you've built}]
One more word for telling, explaining and describing at B1.
\end{grammarlens}

\subsection*{Your turn}
\begin{itemize}
\raggedright
  \item \emph{Say it:} \emph{le stagiaire}
  \item \emph{Say it:} \emph{le stagiaire}, once more
  \item \emph{Say it:} a sentence of your own with \emph{le stagiaire}
  \item \emph{Recall:} say the French for to recruit, then the French for the internship, then say \emph{le stagiaire} again
\end{itemize}

\begin{tcolorbox}[breakable,colback=teal!4,colframe=teal!35!black,title={Before you move on}]
What does le stagiaire mean? (\textquotedblleft{}The intern, the trainee\textquotedblright{}.) Say it once more.
\end{tcolorbox}
\section[l'apprenti]{l'apprenti — the apprentice}

\label{lesson:FR-C357-apprenti}



\begin{quote}
Before the new one: say the French for the internship, then the French for the intern.
\end{quote}

\subsection*{You'll want to know: l'apprenti}
\textbf{l'apprenti} — \textquotedblleft{}the apprentice\textquotedblright{}.

\begin{quote}
Le boulanger a un apprenti de seize ans.
\end{quote}

\begin{quote}
\emph{The baker has a sixteen-year-old apprentice.}
\end{quote}

Feminine \emph{apprentie}.

\begin{cousinweb}
\emph{apprenti} comes from Middle French \emph{apprentis} (EtymDB), built on \emph{apprendre} \textquotedblleft{}to learn\textquotedblright{}, which comes from Latin \emph{apprehendere} \textquotedblleft{}to seize, to grasp\textquotedblright{} (Lewis \& Short).
\end{cousinweb}

\begin{grammarlens}[title={What you've built}]
One more word for telling, explaining and describing at B1.
\end{grammarlens}

\subsection*{Your turn}
\begin{itemize}
\raggedright
  \item \emph{Say it:} \emph{l'apprenti}
  \item \emph{Say it:} \emph{l'apprenti}, once more
  \item \emph{Say it:} a sentence of your own with \emph{l'apprenti}
  \item \emph{Recall:} say the French for the internship, then the French for the intern, then say \emph{l'apprenti} again
\end{itemize}

\begin{tcolorbox}[breakable,colback=teal!4,colframe=teal!35!black,title={Before you move on}]
What does l'apprenti mean? (\textquotedblleft{}The apprentice\textquotedblright{}.) Say it once more.
\end{tcolorbox}
\section[l'apprentissage]{l'apprentissage — the apprenticeship; learning}

\label{lesson:FR-C357-apprentissage}



\begin{quote}
Before the new one: say the French for the intern, then the French for the apprentice.
\end{quote}

\subsection*{You'll want to know: l'apprentissage}
\textbf{l'apprentissage} — \textquotedblleft{}the apprenticeship; learning\textquotedblright{}.

\begin{quote}
L'apprentissage d'une langue prend du temps.
\end{quote}

\begin{quote}
\emph{Learning a language takes time.}
\end{quote}

\emph{Être en apprentissage} — to be an apprentice. Masculine.

\begin{cousinweb}
\emph{apprentissage} is built from \emph{apprenti} and \emph{-age}. \emph{apprenti} rests on \emph{apprendre}, from Latin \emph{apprehendere} \textquotedblleft{}to seize, to grasp\textquotedblright{} (Lewis \& Short).
\end{cousinweb}

\begin{grammarlens}[title={What you've built}]
One more word for telling, explaining and describing at B1.
\end{grammarlens}

\subsection*{Your turn}
\begin{itemize}
\raggedright
  \item \emph{Say it:} \emph{l'apprentissage}
  \item \emph{Say it:} \emph{l'apprentissage}, once more
  \item \emph{Say it:} a sentence of your own with \emph{l'apprentissage}
  \item \emph{Recall:} say the French for the intern, then the French for the apprentice, then say \emph{l'apprentissage} again
\end{itemize}

\begin{tcolorbox}[breakable,colback=teal!4,colframe=teal!35!black,title={Before you move on}]
What does l'apprentissage mean? (\textquotedblleft{}The apprenticeship; learning\textquotedblright{}.) Say it once more.
\end{tcolorbox}
\section[l'employeur]{l'employeur — the employer}

\label{lesson:FR-C357-employeur}



\begin{quote}
Before the new one: say the French for the apprentice, then the French for the apprenticeship.
\end{quote}

\subsection*{You'll want to know: l'employeur}
\textbf{l'employeur} — \textquotedblleft{}the employer\textquotedblright{}.

\begin{quote}
Mon employeur me paie à la fin du mois.
\end{quote}

\begin{quote}
\emph{My employer pays me at the end of the month.}
\end{quote}

The feminine \emph{employeuse} is rare.

\begin{cousinweb}
\emph{employeur} is built from \emph{employer} and \emph{-eur}. \emph{employer} is inherited from Latin \emph{implicāre} \textquotedblleft{}to enfold, to involve\textquotedblright{} (EtymDB; Lewis \& Short).
\end{cousinweb}

\begin{grammarlens}[title={What you've built}]
One more word for telling, explaining and describing at B1.
\end{grammarlens}

\subsection*{Your turn}
\begin{itemize}
\raggedright
  \item \emph{Say it:} \emph{l'employeur}
  \item \emph{Say it:} \emph{l'employeur}, once more
  \item \emph{Say it:} a sentence of your own with \emph{l'employeur}
  \item \emph{Recall:} say the French for the apprentice, then the French for the apprenticeship, then say \emph{l'employeur} again
\end{itemize}

\begin{tcolorbox}[breakable,colback=teal!4,colframe=teal!35!black,title={Before you move on}]
What does l'employeur mean? (\textquotedblleft{}The employer\textquotedblright{}.) Say it once more.
\end{tcolorbox}
\section[(dialogue)]{Retour sur les chapitres 350–357}

\label{lesson:FR-R357-revision-350-357}



\begin{quote}
Say the French for the impact and the repercussion.
\end{quote}

\subsection*{Your turn}
Say these aloud:

\begin{itemize}
\raggedright
  \item the impact, the repercussion, to set off, to arouse, the motive
  \item to react, the reaction, the pretext, to justify, the justification
  \item to prevent, to hinder, to favour, the asset, the brake
  \item the remedy, to remedy, to make worse, to lessen, to make up for
  \item the CV, the candidate, the application, to apply, the hiring
  \item to dismiss, the dismissal, the resignation, to resign, the notice
  \item the job, the unemployment, the unemployed person, to hire, to recruit
  \item the internship, the intern, the apprentice, the apprenticeship, the employer
\end{itemize}

\begin{tcolorbox}[breakable,colback=teal!4,colframe=teal!35!black,title={Before you move on}]
The impact? (\textbf{l'impact}.) The repercussion? (\textbf{la répercussion}.) To set off? (\textbf{déclencher}.) To arouse? (\textbf{susciter}.) The motive? (\textbf{le motif}.) To react? (\textbf{réagir}.) The reaction? (\textbf{la réaction}.) The pretext? (\textbf{le prétexte}.) To justify? (\textbf{justifier}.) The justification? (\textbf{la justification}.) To prevent? (\textbf{empêcher}.) To hinder? (\textbf{entraver}.) To favour? (\textbf{favoriser}.) The asset? (\textbf{l'atout}.) The brake? (\textbf{le frein}.) The remedy? (\textbf{le remède}.) To remedy? (\textbf{remédier}.) To make worse? (\textbf{aggraver}.) To lessen? (\textbf{atténuer}.) To make up for? (\textbf{compenser}.) The CV? (\textbf{le CV}.) The candidate? (\textbf{le candidat}.) The application? (\textbf{la candidature}.) To apply? (\textbf{postuler}.) The hiring? (\textbf{l'embauche}.) To dismiss? (\textbf{licencier}.) The dismissal? (\textbf{le licenciement}.) The resignation? (\textbf{la démission}.) To resign? (\textbf{démissionner}.) The notice? (\textbf{le préavis}.) The job? (\textbf{l'emploi}.) The unemployment? (\textbf{le chômage}.) The unemployed person? (\textbf{le chômeur}.) To hire? (\textbf{embaucher}.) To recruit? (\textbf{recruter}.) The internship? (\textbf{le stage}.) The intern? (\textbf{le stagiaire}.) The apprentice? (\textbf{l'apprenti}.) The apprenticeship? (\textbf{l'apprentissage}.) The employer? (\textbf{l'employeur}.)
\end{tcolorbox}
